\documentclass[11pt]{letter}
\usepackage[margin=2.5cm]{geometry}

\signature{David Alarc\'on\\Universidad Pablo de Olavide, Sevilla, Spain\\dalarub@upo.es}
\address{David Alarc\'on\\Departamento de Econom\'ia, M\'etodos Cuantitativos e Historia Econ\'omica\\Universidad Pablo de Olavide\\Ctra.\ de Utrera, km 1\\41013 Sevilla, Spain}

\begin{document}
\begin{letter}{Editorial Board\\Communications in Mathematical Physics}

\opening{Dear Editors,}

I am submitting the manuscript ``Gap ratio statistics of Riemann zeros:
measurement, mechanism, and the Berry-Keating correction'' for consideration
in \textit{Communications in Mathematical Physics}.

This paper presents the first precision measurement of the rate at which
the gap ratio statistic of Riemann zeta zeros converges to the GUE
prediction, using 21 independent data points up to $\log T = 24$.
We find $\langle r\rangle(T) = 0.59891(13) + 1.245(40)/\log^2 T$,
confirming the Berry-Keating $O(1/\log^2 T)$ prediction quantitatively.

We identify the physical mechanism: Riemann zeros have a narrower
spacing distribution than GUE, contributing $+128\%$ to the correction
coefficient, partially offset by increased anti-correlation ($-29\%$).
This decomposition reproduces 99.5\% of the measured coefficient.

This work is part of a series of companion papers: Paper~B (submitted to
\textit{J.\ Number Theory}) derives the coefficient from first principles
via the Conrey-Snaith formula, and Paper~C (submitted to \textit{Annals
of Mathematics}) proves that if the convergence pattern holds exactly,
the Riemann Hypothesis is true.

The manuscript has not been published elsewhere and is not under
consideration by another journal. All authors have approved the manuscript.

\closing{Sincerely,}

\end{letter}
\end{document}
